\documentclass[11pt]{article}
\usepackage{amsmath,amssymb,amsthm,hyperref}
\newtheorem{proposition}{Proposition}
\newcommand{\per}{\operatorname{per}}
\title{Reflection-root partitions as structured tests of Marcus's conjecture}
\author{Research notes; finite searches only; no priority claim}
\date{6 November 2026}
\begin{document}
\maketitle
\noindent\textbf{Status.} These searches produce no counterexample.
They exhaust specified partitions of fixed root sets, not arbitrary
positive semidefinite matrices or all ranks. The determinant identity
below supplies an exact large permanent without enumerating its terms.

\paragraph{Root sets and their products.}
In $\mathbb R^r$ consider the following lists of vectors, with one vector
for each listed linear form:
\[
 \begin{array}{c|c|c}
 \text{label}&\text{linear forms}&N\\\hline
 A&x_i-x_j\ (i<j)&r(r-1)/2\\
 D&x_i-x_j,\ x_i+x_j\ (i<j)&r(r-1)\\
 B&x_i\ (1\le i\le r),\ x_i-x_j,\ x_i+x_j\ (i<j)&r^2
 \end{array}
\]
The label $A$ denotes the roots in $r$ coordinates, of rank $r-1$.
Let $\mathcal A$ be the $N\times N$ Gram matrix of the chosen list.
For a partition into $m=N/k$ sets of size $k$, form its matrix $G$ of
$k\times k$ block permanents. The value $\per\mathcal A$ is independent
of the partition; only $\per G$ needs to be maximized.

\begin{proposition}
For the three root lists,
\[
 \per\mathcal A=r!\prod_{j=0}^{r-1}d_j!,\qquad
 d_j=j\ (A),\quad d_j=2j\ (D),\quad d_j=2j+1\ (B).
\]
\end{proposition}
\begin{proof}
The product of all listed forms equals, up to an overall sign,
$\det[x_i^{d_j}]_{i=1,\ldots,r;\ j=0,\ldots,r-1}$.
For $A$ this is the Vandermonde formula; for $D$ apply that formula to
$x_i^2$; for $B$ multiply the latter by $\prod_i x_i$.
All $r!$ determinant monomials are distinct because the degrees $d_j$
are distinct. They have coefficients of absolute value one and each has
Fock norm squared $\prod_jd_j!$. Their orthogonality gives the displayed
norm. Gaussian moment expansion equates that norm with $\per\mathcal A$.
\end{proof}

\paragraph{Completed exact searches.}
The standalone program \texttt{marcus-reflection-search.mjs} enumerates
every unlabeled partition in the rows marked ``all.'' It fixes the least
unused root in the next block, then chooses its other $k-1$ members.
The number of leaves is independently checked against
$N!/((k!)^m m!)$. Block entries are cached and their permanents are
calculated by subset dynamic programming. Every intermediate product and
sum is required to stay in the exact integer range of JavaScript numbers.
Thus integer equality is used throughout, not a floating-point comparison
of nearly equal quantities. The maximizing partition in each row is
checked again using the original Gram vectors and an independent BigInt
permanent implementation.
\[
\begin{array}{c|r|r|r|r|r|l}
 \text{type}&r&k&m&\text{partitions}&\per\mathcal A&\max\per G\\\hline
 A&5&2&5&945&34560&4444\\
 A&6&3&5&1401400&24883200&1069134\\
 A&6&5&3&126126&24883200&4746560\\
 D&4&2&6&10395&829440&34960\\
 D&4&3&4&15400&829440&69632\\
 D&4&4&3&5775&829440&111152
\end{array}
\]
All six rows are exhaustive for their stated root set and block size.
The total is $1560041$ partitions. In a separate finite random search of
$30000$ partitions of $B$ at $r=4,k=4,m=4$, the largest observed value
was $5779608$, while $\per\mathcal A=87091200$; no optimality is asserted
for that row. Its reproducible generator seed is $117361$.

For example, the third exhaustive row is reproduced by
\texttt{node marcus-reflection-search.mjs A 6 5 all}; the random row by
\texttt{node marcus-reflection-search.mjs B 4 4 random 30000}.
Each output JSON stores the root vectors and the maximizing observed
partition, as well as exact integer permanent values.

\paragraph{What the experiment changes.}
The reflection construction tests an organized source of extensive
cancellation: a product of many linear forms collapses to a determinant
with only $r!$ monomials. Despite that cancellation, every tested partition
has a substantial positive gap. These particular fixed root systems do
not yield a disproof. Nothing here bounds the effect of arbitrary
deformations, other root systems, or arbitrary Gram matrices. The full
Marcus conjecture and its general equality clause remain unresolved.
\end{document}
